\documentclass[12pt]{article}
\usepackage[utf8]{inputenc}
\usepackage[spanish]{babel}
\usepackage[a4paper, margin=2cm]{geometry}
\usepackage{tikz}
\usetikzlibrary{babel}
\usepackage[most]{tcolorbox}
\usepackage{amsmath}

\renewcommand{\familydefault}{\sfdefault}

% Definición de colores
\definecolor{medblue}{RGB}{43, 108, 176}
\definecolor{medteal}{RGB}{49, 151, 149}
\definecolor{medlight}{RGB}{235, 248, 255}
\definecolor{meddark}{RGB}{26, 32, 44}
\definecolor{aortared}{RGB}{239, 68, 68}
\definecolor{bloodred}{RGB}{185, 28, 28}
\definecolor{valveyellow}{RGB}{253, 224, 71}

\begin{document}
\pagestyle{empty}

\begin{tcolorbox}[
    enhanced, 
    colback=medlight!30, 
    colframe=medblue, 
    title={\Large \textbf{Cinemática: Aceleración Sanguínea}}, 
    halign title=center, 
    fonttitle=\bfseries, 
    arc=4mm, 
    boxrule=1.5pt,
    shadow={2mm}{-2mm}{0mm}{black!30!white}
]

    \vspace{0.2cm}
    \begin{tcolorbox}[colback=white, colframe=medteal, arc=2mm, boxrule=1.2pt, title=\textbf{Caso Clínico / Problema}]
        Durante la fase de sístole, la sangre es bombeada desde el ventrículo izquierdo hacia la aorta, acelerando desde el reposo hasta una velocidad de \textbf{0.4 m/s} en un tiempo de \textbf{0.1 s}. Calcule la aceleración media de la sangre en este proceso.
    \end{tcolorbox}

    \vspace{0.4cm}
    
    \begin{center}
    \begin{tikzpicture}[scale=1]
        
        % Vaso sanguíneo (Aorta) simulando expansión
        \filldraw[aortared!20, draw=aortared!80!black, ultra thick] (0, 0) .. controls (2, 0.5) and (6, 0.5) .. (8, 0) -- (8, 2.5) .. controls (6, 2) and (2, 2) .. (0, 2.5) -- cycle;
        
        % Válvula aórtica (Izquierda)
        \filldraw[valveyellow, draw=meddark, thick] (0, 0) .. controls (0.6, 0.6) and (0.6, 1.9) .. (0, 2.5);
        \node[font=\bfseries, text=meddark, align=center, rotate=90] at (-0.4, 1.25) {Ventrículo\\Izquierdo};

        % Volumen de sangre t=0 (Inicio)
        \filldraw[bloodred, draw=meddark, thick, rounded corners=3pt] (1.0, 0.8) rectangle (1.8, 1.7);
        \node[font=\bfseries, text=bloodred, below] at (1.4, -0.2) {$v_i = 0$};
        \node[font=\bfseries, text=meddark, above, fill=white, inner sep=2pt, rounded corners=2pt] at (1.4, 2) {$t = 0 \text{ s}$};

        % Volumen de sangre t=0.1 (Final)
        \filldraw[bloodred, draw=meddark, thick, rounded corners=3pt] (5.5, 0.8) rectangle (6.3, 1.7);
        \draw[->, ultra thick, medteal] (6.4, 1.25) -- (8.5, 1.25) node[right, font=\bfseries] {$v_f = 0.4 \text{ m/s}$};
        \node[font=\bfseries, text=meddark, above, fill=white, inner sep=2pt, rounded corners=2pt] at (5.9, 2) {$t = 0.1 \text{ s}$};

        % Flecha global de aceleración
        \draw[->, ultra thick, meddark, dashed] (2.2, 2.8) -- (5.0, 2.8) node[midway, above, font=\Large\bfseries] {Aceleración $\vec{a}$};
        
        % Títulos de referencia
        \node[font=\bfseries, text=aortared!80!black] at (3.5, -0.6) {Arteria Aorta};

    \end{tikzpicture}
    \end{center}

    \vspace{0.4cm}
    \begin{tcolorbox}[colback=medteal!10, colframe=medteal, arc=2mm, boxrule=1.2pt, title=\textbf{Desarrollo Matemático y Solución}]
        La aceleración media ($a$) se define como el cambio de velocidad dividido por el tiempo transcurrido en realizar dicho cambio. Utilizando la ecuación cinemática básica:
        \begin{align*}
            a &= \frac{v_f - v_i}{t} \\[1ex]
            a &= \frac{0.4 \text{ m/s} - 0 \text{ m/s}}{0.1 \text{ s}} = \frac{0.4}{0.1} \text{ m/s}^2
        \end{align*}
        \textbf{Resultado:} La aceleración media de la sangre al ingresar a la aorta es de \textbf{4 m/s$^2$}.
    \end{tcolorbox}
    
\end{tcolorbox}

\end{document}